
%\clearpage
\section{\MakeUppercase{Fundamentação Legal}}
\label{FUNDAMENTAÇÃO LEGAL}

A elaboração e a atualização deste Projeto Pedagógico de Curso observam o conjunto de normativas vigentes no âmbito nacional e institucional que regulamentam os cursos superiores, em especial aqueles classificados como bacharelados. Tais diretrizes asseguram a coerência acadêmica, a consistência metodológica e a aderência do curso às políticas públicas de Educação Superior no Brasil.

Ademais, a presente reformulação do PPC contempla as determinações recentes referentes à organização curricular, às diretrizes de formação profissional, às atividades de extensão e ao processo de avaliação interna e externa do curso, seguindo as orientações dos órgãos reguladores nacionais e das instâncias normativas do IFCE.

\subsection{Normativas Nacionais}

\begin{itemize}
	\item Lei Nº 9.394, de 20 de dezembro de 1996, que estabelece as Diretrizes e Bases da Educação Nacional;
	
	\item Resolução CNE/CES 02/2019 – Institui Diretrizes Curriculares Nacionais do Curso de graduação em Engenharia;
	
	\item Parecer CNE/CES n° 01/2019 que trata das Diretrizes Curriculares Nacionais do Curso de Graduação em Engenharia;	
	
	\item Lei nº 5.194/66 – Regula o exercício das profissões de Engenheiro, Arquiteto e Engenheiro-Agrônomo, e dá outras providências;
	
	\item Resolução CNE/CES nº 02/2007 e Parecer CNE/CES nº. 08/2007 – Dispõem sobre carga horária mínima e procedimentos relativos à integralização e duração dos cursos de graduação, bacharelados, na modalidade presencial, bem como estabelecem que os estágios e atividades complementares dos cursos de graduação, bacharelados, na modalidade presencial,
	não deverão exceder a 20\% (vinte por cento) da carga horária total do curso;
	
	\item Lei Nº 11.788, de 25 de setembro de 2008, que dispõe sobre o estágio de estudantes, altera a redação do art. 428 da Consolidação das Leis do Trabalho – CLT e dá outras providências;
	
	
	\item Decreto nº 5.626, de 22/12/2005, que regulamenta a Lei nº 10.436, de 24/04/2002, que dispõe sobre a Língua Brasileira de Sinais – LIBRAS como disciplina curricular obrigatória nos cursos de formação de professores para o exercício do magistério, em nível médio e superior, e nos cursos de Fonoaudiologia, bem como disciplina curricular optativa nos
	demais cursos de educação superior e na educação profissional (BRASIL, 2005);
	
	\item Lei N°. 13.146, de 06 de 2015, que institui a Lei Brasileira de Inclusão da Pessoa com Deficiência (Estatuto da Pessoa com Deficiência);
	
	\item  Resolução CNE/CP nº 2, de 15 de junho de 2012, que estabelece as Diretrizes Curriculares Nacionais para a Educação Ambiental;
	
	\item Resolução CNE/CP nº 1, de 30 de maio de 2012. Estabelece as Diretrizes Nacionais para a Educação em Direitos Humanos;

	\item Resolução CNE/CP nº 1, de 17 de junho de 2004. Institui Diretrizes Curriculares Nacionais para a Educação das Relações Étnico-Raciais e para o Ensino de História e Cultura 	Afro-Brasileira e Africana;
	
	\item Lei Nº 13.005, de 25/6/2014 - Aprova o Plano Nacional de Educação - PNE e dá outras providências;
	
	\item Portaria Nº 2.117, DE 6 DE DEZEMBRO DE 2019 Dispõe sobre a oferta de carga horária na modalidade de Ensino a Distância - EaD em cursos de graduação presenciais ofertados por Instituições de Educação Superior - IES pertencentes ao Sistema Federal de Ensino;
	
	\item Resolução nº 7, de 18/12/2018 - Estabelece as Diretrizes para a Extensão na
	Educação Superior Brasileira e regimenta o disposto na Meta 12.7 da Lei nº13.005/2014, que aprova o Plano Nacional de Educação - PNE 2014-2024 e dá outras providências;
	
	\item Lei Nº 11788, de 25 de setembro de 2008 - Dispõe sobre o estágio de estudantes e altera a redação do art. 428 da Consolidação das Leis do Trabalho – CLT.
	
\end{itemize}
%
\subsection{Normativas Institucionais}
%
\begin{itemize}
	\item Regulamento da Organização Didática do IFCE (ROD);
	\item Plano de Desenvolvimento Institucional (PDI);
	\item Projeto Pedagógico Institucional (PPI);
	\item Resolução N° 100, de 27 de setembro de 2017 - CONSUP e suas alterações: estabelece os procedimentos para criação, suspensão e extinção de cursos no IFCE;
	\item Portaria Nº 176/GABR/REITORIA, de 10 de maio de 2021: atualização do Perfil Docente do IFCE;
	\item Resolução Consup nº 028, de 08 de agosto de 2014, que dispõe sobre o Manual de Estágio do IFCE;
	
	\item Resolução Nº 39, de 22 de agosto de 2016: Aprova alteração na redação das artigos 7°, 9°(incluindo quadro I), 10°(incluindo quadro 2) e 12°(incluindo quadro 3) e anexos I,II e III da Regulamentação das Atividades Docentes(RAD) do IFCE.
	
	\item Resolução Nº 63, de 28 de maio de 2018: Aprova a Regulamentação das Atividades Docentes (RAD) do Instituto Federal de Educação, Ciência e Tecnologia do Ceará -
IFCE.
	\item Resolução CONSUP Nº 004/2015, que regulamenta o Núcleo Docente Estruturante (NDE) no IFCE;
	
	\item Guia de Curricularização das Atividades de Extensão nos cursos técnicos, de graduação e pós-graduação do IFCE (\citeonline{GuiaCurricExt2023});
	
	\item Resolução CONSUP Nº 75/2018, que estabelece normas de funcionamento dos colegiados de curso;
	
	\item Resolução CONSUP/IFCE nº 259, de 8 de janeiro de 2025 (\citeonline{res259}), que institui o \textit{Manual de Normalização de Projetos Pedagógicos dos Cursos de Engenharia do IFCE}.

\end{itemize}

